\section{Graph differences and even subsets}
For a nonempty set $\Omega$, let $\F_2^{(\Omega)}$ be the vector space
of finitely supported functions on $\Omega$, with coordinate vectors
$e_v$. Its augmentation module is
\[
 E(\Omega)=\left\{z\in\F_2^{(\Omega)}:
                 \sum_{v\in\Omega}z_v=0\right\}.
\]
For a chosen vertex $o$, the vectors $e_o+e_v$, $v\ne o$, are a basis.
Indeed a vector of even total sum is the sum of $e_o+e_v$ over its
nonzero coordinates away from $o$; independence follows by inspecting
those coordinates. Thus, if $|\Omega|=m<\infty$, the module has order
$2^{m-1}$.

\begin{theorem}[Presentation by graph differences]\label{UC.graphmodule}
Let $Q$ be a connected simple graph with vertex set $\Omega$. Define
$B(Q)$ by one generator $z_e$ for every edge, with relations
\[
 z_e^2=1,\qquad [z_e,z_f]=1,
 \qquad z_{e_1}\cdots z_{e_k}=1
\]
for every finite closed walk with successive edges $e_1,\ldots,e_k$.
Then $z_{\{u,v\}}\mapsto e_u+e_v$ defines an isomorphism
$B(Q)\cong E(\Omega)$. It commutes with every graph automorphism.
\end{theorem}
\begin{proof}
The edge vectors have order dividing two and commute. Their sum along a
closed walk is zero because each successive vertex occurs twice. Hence
the indicated homomorphism exists.

Choose a base vertex $o$. For each $v$, choose a finite path from $o$ to
$v$ and let $t_v$ be the product of its edge generators. The product is
independent of the chosen path: concatenate one path with the reverse
of the other, and apply the closed-walk relation. In particular $t_o=1$.
Following a path to $u$, then the edge $\{u,v\}$, and finally a reverse
path from $v$ to $o$ gives $z_{\{u,v\}}=t_ut_v$.

The elements $t_v$ commute and have square one. The basis description
above therefore defines a homomorphism $E(\Omega)\to B(Q)$ sending
$e_o+e_v$ to $t_v$. The two composites fix every edge generator and
every basis vector, respectively, so they are identity maps.
The original edge map, unlike the chosen inverse formula, uses no
base vertex or paths. It consequently commutes with every graph
automorphism.
\end{proof}

For a finite graph one may use a spanning tree and only the fundamental
cycles determined by the remaining edges. The same proof applies:
the unique tree paths define $t_v$, and each fundamental-cycle relation
gives $z_{\{u,v\}}=t_ut_v$ for its extra edge. This also proves directly
that every closed-walk relation follows from that finite list.

\begin{corollary}[Generation inside a permutation extension]\label{UC.edgegeneration}
Let $H$ act by graph automorphisms on a connected graph $Q$. In
$E(\Omega)\rtimes H$, suppose a subgroup contains $\{0\}\rtimes H$
and translations $e_u+e_v$ for a set of edges whose $H$-orbits form a
connected spanning subgraph. Then it is the whole semidirect product.
\end{corollary}
\begin{proof}
Conjugation by $(0,h)$ carries the translation $e_u+e_v$ to
$e_{hu}+e_{hv}$. The subgroup therefore contains every edge difference
in the indicated spanning subgraph. Theorem~\ref{UC.graphmodule} says
that these generate $E(\Omega)$, proving the assertion.
\end{proof}

\section{Recovering a semidirect product from its translations}
Let $H$ act on a nonempty set $\Omega$, and fix $o\in\Omega$.
We use left actions and multiplication
\[
 (v,g)(w,h)=(v+gw,gh)
       \quad\hbox{in }E(\Omega)\rtimes H.
\]
The following statement does not require either set or group to be
finite.

\begin{theorem}[Recognition by differences from a base vertex]\label{UC.recognition}
Let $P$ be a group, let $j:H\to P$ be a homomorphism, and let
$(t_v)_{v\in\Omega}$ be commuting elements of $P$ satisfying
$t_v^2=1$ and $t_o=1$. Suppose
\begin{equation}\label{UC.covariance}
 j(h)t_vj(h)^{-1}=t_{ho}t_{hv}
        \qquad(h\in H,\ v\in\Omega).
\end{equation}
There is a unique homomorphism
\[
 \Psi:E(\Omega)\rtimes H\longrightarrow P,
 \qquad \Psi(e_o+e_v,1)=t_v,
 \qquad \Psi(0,h)=j(h).
\]
It is surjective if these elements generate $P$. Under that generation
hypothesis, if a homomorphism $\Phi:P\to E(\Omega)\rtimes H$ has
the displayed inverse values on these elements, then $\Phi$ and
$\Psi$ are mutually inverse.
\end{theorem}
\begin{proof}
The commuting involutions define a homomorphism
$T:E(\Omega)\to P$ by $T(e_o+e_v)=t_v$ on the stated basis.
The action of $h$ on this vector is
\[
 h(e_o+e_v)=(e_o+e_{ho})+(e_o+e_{hv}).
\]
Consequently~\eqref{UC.covariance} implies
$T(hw)=j(h)T(w)j(h)^{-1}$ first on every basis vector and then on
every vector. Define $\Psi(w,h)=T(w)j(h)$. For two such pairs,
\[
 T(v)j(g)T(w)j(h)=T(v)T(gw)j(gh)=T(v+gw)j(gh),
\]
which is exactly the required homomorphism identity. Uniqueness and
surjectivity follow from the generating sets. Under the final
hypothesis each composite fixes its respective generating set, and is
therefore the identity.
\end{proof}

\begin{lemma}[A twist exchanging two complements]\label{UC.twist}
The function $c(h)=e_o+he_o$ satisfies
$c(gh)=c(g)+gc(h)$. Thus
\[
 \delta(v,h)=(v+c(h),h)
\]
is an involutory automorphism of $E(\Omega)\rtimes H$. The standard
complement $A=\{(0,h):h\in H\}$ and its image satisfy
\[
 A\cap\delta(A)=\{(0,h):h\in H_o\},
\]
where $H_o$ is the stabilizer of $o$.
\end{lemma}
\begin{proof}
In characteristic two,
$e_o+ghe_o=(e_o+ge_o)+g(e_o+he_o)$, proving the cocycle identity.
Substitution into the multiplication rule proves that $\delta$ is a
homomorphism; applying it twice adds $c(h)$ twice, so gives the identity.
Finally $(0,h)$ belongs to its image of $A$ exactly when
$c(h)=0$, or equivalently $ho=o$.
\end{proof}

\section{Cyclic displacement relations}
The next identity is the displacement argument in Li--Zhang's proof of
their Theorem~3.4. Keeping integer indices exposes the action needed in
Theorem~\ref{UC.recognition} and avoids a separate order argument.

\begin{lemma}\label{UC.displacements}
Let $x,y$ belong to a group and put $t_k=y^kx^k$ for $k\in\Z$.
If every $t_k$ has square one, then all $t_k$ commute and
\[
 x^{-1}t_kx=t_1t_{k+1},\qquad
 xt_kx^{-1}=t_{-1}t_{k-1}.
\]
If an involution $b$ inverts $x$ and $y$, then $bt_kb=t_{-k}$.
If $x^n=y^n=1$, it suffices to require $t_k^2=1$ for
$1\le k<n$.
\end{lemma}
\begin{proof}
From $t_k=t_k^{-1}$ we obtain
$y^kx^k=x^{-k}y^{-k}$ and, on conjugating by $x^k$,
$x^ky^k=y^{-k}x^{-k}$. Hence
\[
 \begin{aligned}
 t_it_j
 &=x^{-i}y^{j-i}x^j
 =x^{-j}(x^{j-i}y^{j-i})x^j\\
 &=x^{-j}y^{i-j}x^i=t_jt_i.
 \end{aligned}
\]
Direct multiplication gives
$x^{-1}t_kx=t_1^{-1}t_{k+1}=t_1t_{k+1}$.
For $k=-1$ this yields $xt_1x^{-1}=t_{-1}$.
Conjugate the first formula with index $k-1$ by $x$ and use this
identity to obtain the second formula. Conjugation by $b$ gives
$bt_kb=y^{-k}x^{-k}=t_{-k}$. Finally the two period relations give
$t_{k+n}=t_k$ for every integer $k$, while $t_0=1$; this proves the
finite-relation assertion.
\end{proof}

\begin{corollary}\label{UC.cyclicrecognition}
Suppose $n\ge3$ and $P=\gen{x,b,y}$ satisfies
\[
 x^n=y^n=b^2=1,\quad bxb=x^{-1},\quad byb=y^{-1},
 \quad(y^kx^k)^2=1\quad(1\le k<n).
\]
Let $H=D_{2n}$ act on $\Z/n\Z$ by $a(i)=i+1$ and $b(i)=-i$.
If there is a homomorphism $\Phi:P\to E(\Z/n\Z)\rtimes H$ with
\[
 x\longmapsto(0,a),\quad b\longmapsto(0,b),\quad
 y\longmapsto(e_0+e_{-1},a^{-1}),
\]
then it is an isomorphism. Its inverse sends
$(e_0+e_i,1)$ to $t_{-i}$ and $(0,a),(0,b)$ to $x,b$.
\end{corollary}
\begin{proof}
The dihedral relations define $j:H\to P$, with $j(a)=x$ and
$j(b)=b$. Set $T_i=t_{-i}$; the period relations make this independent
of the integer representative of $i$. Lemma~\ref{UC.displacements}
gives commuting involutions and
\[
 xT_ix^{-1}=T_1T_{i+1},\qquad bT_ib=T_{-i}.
\]
These are exactly~\eqref{UC.covariance} for $a,b$, with $o=0$.
Equivariance for products and inverses follows by applying these
identities successively; thus it holds for all $h\in H$.
The elements $j(H),T_i$ generate $P$, since $y=t_1x^{-1}$.
Moreover $\Phi(t_k)=(e_0+e_{-k},1)$ by direct multiplication, or by
the cocycle identity in Lemma~\ref{UC.twist}. The final assertion of
Theorem~\ref{UC.recognition} now gives the inverse and the isomorphism.
\end{proof}
